\chapter{Этаж~1. Чуть выше: облако и \texorpdfstring{$\varepsilon$}{epsilon}-оболочки}
\label{ch:floor1}

\section{Радиус этажа}
\label{sec:floor1-scale}

Этаж~1 --- не одна ячейка, а \emph{составной узел} на каузальной звезде FCC.
Длина задаётся степенями числа соседей $N_{12}=12$:
\[
  N_{12}^{1}=12,\quad
  N_{12}^{2}=144,\quad
  N_{12}^{3}=1728,\quad
  N_{12}^{4}\approx 2.1\times 10^{4},\quad
  N_{12}^{5}\approx 2.5\times 10^{5}.
\]
Это окно $10^{3}\ldots 10^{5}$ планковских шагов~$\hl$ ---
\emph{до} комптоновской и боровской шкал ($10^{22}\ldots 10^{24}\cdot\hl$,
\cref{prop:ladder-scale-steps}).

Здесь живут оболочки соседей вокруг ядра планкона и облако материи~$\rho_\Theta$,
аналог электронной плотности на атомном языке, но на~$\Mlayer$.

\section{Три слоя плотности}
\label{sec:floor1-density-layers}

На одном узле различают три объекта, которые на~$T$ часто смешивают:

\begin{table}[htbp]
  \centering
  \caption{Слои плотности вокруг планкона.}
  \label{tab:floor1-density}
  \begin{tabular}{@{}lll@{}}
    \toprule
    поле & смысл на~$\Mlayer$ & аналог на~$T$ \\
    \midrule
    $|z|^2$, $\rho_E$ & энергоплотность \emph{среды}; вакуум кипит везде & не $\rho_e$ \\
    $b\in\{0,1\}$ на ядре & занятость ячейки: дырка или планкон & жёсткий скелет узла \\
    $\rho_\Theta(x)$ & облако вокруг ядра: где $|\Delta\varphi|$, $|\zeta|\geq\Delta\varphi_{\min}$ & $\rho_e(r)$ \\
    \bottomrule
  \end{tabular}
\end{table}

Гейзенберг на голономии ($\Delta\varphi\geq\Delta\varphi_{\min}=\tfrac12$ рад,
\cref{lem:51a}) не позволяет сжать фазовый дефект в бесконечно острый угол.
Поэтому планкон не может быть только битом «ячейка занята'':
вокруг ядра остаётся \emph{неустранимый ореол}.
Модуль~$K_P$ (anti-smear) держит ядро; Гейзенберг держит хвост.

На~$\Mlayer$ облако задаётся порогом, без подгоняемой гладкой функции:
\[
  \rho_\Theta(x)=\mathbf{1}\big[|\Delta\varphi_N(x)|\geq\Delta\varphi_{\min}\big].
\]
Гладкость появляется только после осреднения~\cref{def:soft-readout} на~$T$
(binomial coarse-grain).

\section{Минимальный ореол: \texorpdfstring{$R_{\mathrm{dress}}=\hl$}{R=dl}}
\label{sec:floor1-dressing}

\begin{proposition}[Радиус одевания]
\label{prop:floor1-dress-radius}
У легчайшего заряженного состояния $Q=\pm1$ минимальный пространственный хвост
облака~$\rho_\Theta$ покрывает всю $\varepsilon$-звезду первых соседей:
\[
  R_{\mathrm{dress}}=R_{\min}=1\cdot\hl,\qquad N_{\mathrm{star}}=N_{12}.
\]
\end{proposition}

\begin{proof}[Идея доказательства]
На звезде из $N_{12}$ связей фазовый шаг за оборот
$\Delta\varphi_{\mathrm{ring}}=2\pi/N_{12}$.
Из $N_\varphi=\lceil 2\pi/\Delta\varphi_{\min}\rceil=13$ следует
$2\pi/N_{12}>\Delta\varphi_{\min}$: вся $\varepsilon$-окрестность лежит выше
Гейзенберг-пола.
Локальный закон~$g$ действует только на $N(x)$ за один такт (A1);
вторая координационная сфера за один~$\htick$ недоступна.
Возбуждённое multi-shell облако --- не основное состояние того же $Q$
и спускается к минимальному ореолу (логика классов C3 vs C2 ниже).
\end{proof}

Оболочки на этаже~1 --- это моды~$\rho_\Theta$ на координационных сферах,
не орбитали $r<\ell_P$ внутри пикселя (\cref{ch:floor0}).

\section{Кто устойчив на этаже~1}
\label{sec:floor1-census}

На радиусах $N_{12}^{3}\ldots N_{12}^{4}$ рассматривают классы объектов с $B=0$
(нейтральный сектор по магнитному/топологическому бюджету узла).
Критерий устойчивости: нет пути вниз по~$g$ с теми же инвариантами и меньшей энергией.

\begin{table}[htbp]
  \centering
  \caption{Классы на лептонном окне этажа~1.}
  \label{tab:floor1-census}
  \begin{tabular}{@{}llp{0.42\linewidth}@{}}
    \toprule
    класс & устойчив? & комментарий \\
    \midrule
    C0 вакуум / A5-фон & нет & не объект, среда \\
    C1 свободный $\gamma$ & да & уже вакуум; не blob радиуса $N_{12}^{k}$ \\
    C2 легчайший $Q=\pm1$ + ореол & \textbf{да} & ядро этажа~0 + $\rho_\Theta$; \cref{prop:floor1-dress-radius} \\
    C3 возбуждённый тот же $Q$ & нет & канал $C3\to C2+\gamma$ существует \\
    C4 $Q=0$ мультиклеточный blob & нет & без топозаряда $\to$ пре-резонанс, не материя \\
    C5 пара $\pm$ на этом $R$ & нет & аннигиляция \\
    C6 нейтральный $L$-сектор & возможно & вне полосы $N_{12}^{3\ldots 4}$ \\
    \bottomrule
  \end{tabular}
\end{table}

\emph{Пре-резонанс} в окне $N_{12}^{3}\ldots N_{12}^{5}$ --- это именно неустойчивые
мультиклеточные конфигурации C4, а не «ещё одна частица''.
Единственная устойчивая \emph{материя} $B=0$ на этом радиусе --- одетый легчайший
$Q=\pm1$ (класс C2).

\section{Связь с электроном}
\label{sec:floor1-electron-bridge}

Лабораторный электрон --- не объект этажа~0 и не весь этаж~1 целиком.
Схема:
\begin{enumerate}[noitemsep]
\item этаж~0: планкон ($b=1$, внутренний спектр, \cref{ch:floor0});
\item этаж~1: неустранимый ореол~$\rho_\Theta$ на $\varepsilon$-звезде;
\item дальше: составной узел (этаж~2), IR-лестница, атом на~$T$
  (\cref{ch:floor2,ch:si-sm}).
\end{enumerate}

Масса $m_e$, заряд $e$ и комптоновская длина --- на оси $k\approx 8$,
не на $N_{12}^{m}$.
Путать $R_{\mathrm{dress}}=\hl$ с $\bar\lambda_C$ или $a_0$ --- частая ошибка;
\cref{ch:floors-preface}.

\section{Открыто на этаже~1}
\label{sec:floor1-open}

\begin{itemize}[noitemsep]
\item полный каталог мод~$\rho_\Theta$ на 2-й и 3-й координационных сферах;
\item правила отбора между оболочками (схема есть, sim нет);
\item время жизни возбуждённого C3 в заполненной A5-ванне (не continuum-$\Gamma$);
\item отношение $m_\mu/m_e$ --- \emph{не} размер этажа~1, отдельный leaf.
\end{itemize}

Следующий шаг вверх --- \cref{ch:floor2}: конфайнмент и кварковые фокусы
на радиусе $\sim 10^{15}\cdot\hl$.
